\chapter*{K}
\label{chap:k}

\term{K-theory (Topological)}
A cohomology-like theory that assigns a ring $K(X)$ to a topological space $X$, constructed from the Grothendieck group of vector bundles over $X$. It is a powerful tool for studying the fact that some vector bundles are not trivial.

\term{K-theory (Algebraic)}
A sequence of functors $K_n$ from rings to abelian groups. $K_0(R)$ is the Grothendieck group of finitely generated projective modules over $R$. Algebraic K-theory bridges ring theory, topology, and number theory.

\term{Kähler Manifold}
A complex manifold equipped with a Hermitian metric whose associated 2-form $\omega$ is closed, so that $[\omega]$ defines a Kähler class in $H^2(X, \mathbb{R})$. The condition $d\omega = 0$ makes the Levi--Civita connection and the Chern connection coincide, and it implies the $\partial\bar{\partial}$-lemma, the Hodge decomposition into Dolbeault cohomology groups, and vanishing theorems such as Kodaira's. Complex projective space and smooth projective varieties are Kähler. Kähler geometry is central to complex and algebraic geometry, mirror symmetry, and the study of Calabi--Yau manifolds.

\term{Kantorovich Inequality}
The inequality $(\sum a_i x_i)(\sum a_i / x_i) \leq \frac{(M+m)^2}{4Mm}$ for positive $a_i$ and $m \leq x_i \leq M$, used in optimization and numerical analysis.

\term{Karush--Kuhn--Tucker (KKT) Conditions}
First-order necessary conditions for a solution in nonlinear programming to be optimal, provided that some regularity conditions are met. They generalize the method of Lagrange multipliers to inequality constraints.

\term{Kazhdan--Lusztig Polynomials}
Polynomials associated with Coxeter groups used to describe the singularities of Schubert varieties. They play a central role in the representation theory of Lie algebras and the Hecke algebra.

\term{Keller's Conjecture}
The conjecture that any tiling of $\mathbb{R}^n$ by equal hypercubes has a pair of cubes meeting full-face-to-full-face; true for $n \leq 6$ and disproved in dimension 10 or more.

\term{Kepler Conjecture}
The theorem that the densest packing of equal spheres in three dimensions is the face-centred cubic lattice, with density $\pi/\sqrt{18}$.

\term{Kernel (Linear Algebra)}
The set of all vectors $v$ such that $Tv = 0$ for a linear operator $T$. Also known as the null space. The dimension of the kernel is the nullity of the operator.

\term{Kernel (Group Theory)}
The set of elements in the domain of a group homomorphism that are mapped to the identity element of the codomain. The kernel is always a normal subgroup of the domain.

\term{Kernel (Integral Equation)}
The function $K(x,y)$ in an integral operator $(Tf)(x) = \int K(x,y)f(y)\,dy$. The properties of the kernel (e.g., symmetry, continuity) determine the properties of the operator.

\term{Kernel (Convolution)}
A function used in convolution to extract specific features from a signal or image. In the context of the Fourier transform, convolution in the time domain corresponds to multiplication in the frequency domain.

\term{Kernel Method}
The technique of implicitly mapping data into a high-dimensional feature space via a positive definite kernel, used in support vector machines.

\term{Killing Form}
A symmetric bilinear form on a Lie algebra $\mathfrak{g}$ defined by $B(X, Y) = \operatorname{tr}(\operatorname{ad}_X \operatorname{ad}_Y)$. A Lie algebra is semisimple iff its Killing form is non-degenerate.

\term{Killing Vector Field}
A vector field $X$ on a Riemannian manifold such that the Lie derivative of the metric $g$ along $X$ vanishes: $\mathcal{L}_X g = 0$. These fields generate isometries of the manifold.

\term{Kirchhoff's Laws}
Conservation laws for electric circuits: the current law (sum of currents into a node is zero) and the voltage law (sum of potential differences around a loop is zero).

\term{Kirchhoff's Theorem}
The formula for the number of spanning trees of a graph as any cofactor of its Laplacian matrix.

\term{Klee's Measure Problem}
The problem of computing the volume of the union of $n$ axis-parallel hyperrectangles in $d$-dimensional space. The most efficient algorithms have complexity $O(n \log n)$ for $d=2$ and $O(n^{d/2})$ for larger $d$.

\term{Klein Bottle}
A non-orientable surface of genus 1, formed by gluing the opposite edges of a rectangle with one pair reversed. It cannot be embedded in $\mathbb{R}^3$ without self-intersection, though it can be embedded in $\mathbb{R}^4$.

\term{Klein Four-Group}
The group $\mathbb{Z}/2 \times \mathbb{Z}/2$, the smallest non-cyclic group, with four elements each of order 2.

\term{Klein Quadric}
The Grassmannian $\operatorname{Gr}(2, 4)$ of lines in $\mathbb{P}^3$, which can be embedded as a smooth quadric hypersurface in $\mathbb{P}^5$. It is central to the study of line geometry.

\term{Kneser Graph}
The graph $KG(n,k)$ whose vertices are the $k$-subsets of an $n$-set, adjacent when they are disjoint; used in Ramsey theory.

\term{Knot Theory}
The study of mathematical knots, which are embeddings of the circle $S^1$ into $\mathbb{R}^3$ or $S^3$. Two knots are equivalent if one can be deformed into the other without cutting.

\term{Knot Invariants}
Quantities (numbers, polynomials, or groups) assigned to a knot that remain unchanged under ambient isotopy. Examples include the Jones polynomial, the Alexander polynomial, and the knot group $\pi_1(S^3 \setminus K)$.

\term{Knot Sum}
The connected sum of two knots $K_1$ and $K_2$, performed by removing a small arc from each and joining the endpoints. This operation makes the set of knots into a commutative monoid.

\term{Knuth's Algorithm}
Algorithms of Donald Knuth for combinatorial generation and sorting, notably his efficient permutation generation algorithms.

\term{Konigsberg Bridges}
The historical problem, solved by Euler in 1736, of whether one can walk across all seven bridges of Königsberg exactly once, founding graph theory.

\term{König's Theorem}
In a bipartite graph, the size of a maximum matching equals the size of a minimum vertex cover. It is equivalent to Hall's marriage theorem and follows from the max-flow min-cut theorem.

\term{K-tuple}
An ordered list of $k$ elements. The set of all $k$-tuples of a set $S$ is the Cartesian product $S^k$.

\term{Koopman Operator}
An infinite-dimensional linear operator that describes the evolution of functions (observables) on a state space under a dynamical system, mapping $f \to f \circ T$.

\term{Kolmogorov Complexity}
The length of the shortest computer program that produces a given string as output. It provides a formal measure of the algorithmic randomness of a string.

\term{Kolmogorov--Savage Divergence}
Also known as Kullback--Leibler (KL) divergence. A measure of how one probability distribution $P$ differs from a second, reference probability distribution $Q$.

\term{Kolmogorov Continuity Theorem}
A criterion for a stochastic process to have a continuous modification. It requires the increments to be bounded by a power of the time difference: $\mathbb{E}[|X_t - X_s|^\alpha] \leq C|t-s|^{1+\beta}$.

\term{Kolmogorov--Smirnov Test}
A non-parametric test used to determine if a sample comes from a specific probability distribution by comparing the empirical distribution function to the theoretical one.

\term{Kolmogorov Three-Series Theorem}
A series of independent random variables $\sum X_n$ converges almost surely if and only if three series converge: $\sum P(|X_n| > 1)$, the series of the means of the truncated variables $X_n 1_{\{|X_n| \leq 1\}}$, and the series of their variances.

\term{K-nearest neighbors (KNN)}
A non-parametric method used for classification and regression, where the label of a point is determined by the majority vote of its $k$ closest neighbors in the feature space.

\term{K-means Clustering}
An iterative algorithm that partitions a dataset into $k$ clusters by minimizing the variance within each cluster (the sum of squared distances to the centroid).

\term{Koszul Complex}
A chain complex used in commutative algebra to study the regularity of a sequence of elements in a ring. It is the primary tool for defining and computing the depth of a module.

\term{Koszul Duality}
A duality between certain categories of modules over quadratic algebras. It relates the cohomology of a space to the algebra of its loop space.

\term{Krull Dimension}
The supremum of the lengths of all chains of prime ideals in a commutative ring. For a polynomial ring $k[x_1, \ldots, x_n]$, the Krull dimension is $n$.

\term{Krull's Intersection Theorem}
In a Noetherian ring, if $I$ is an ideal and $M$ is a finitely generated module, the intersection $\bigcap_{n=1}^\infty I^n M$ consists of elements $m \in M$ such that $(1-a)m = 0$ for some $a \in I$.

\term{Krull--Schmidt Theorem}
The theorem that modules of finite length decompose uniquely into indecomposable direct summands.

\term{Kronecker Delta}
A function of two variables $\delta_{ij}$ that is $1$ if $i=j$ and $0$ otherwise. It is the identity element for the Hadamard product and is used to represent the identity matrix.

\term{Kronecker Product}
An operation on two matrices $A$ and $B$ resulting in a block matrix $A \otimes B$ where each element $a_{ij}$ of $A$ is multiplied by the entire matrix $B$.

\term{Kronecker Lemma}
A result in real analysis stating that if $\sum a_n$ is a convergent series of real numbers and $0 < b_n \uparrow \infty$, then $\frac{1}{b_n} \sum_{k=1}^n a_k b_k \to 0$.

\term{Kronecker--Weber Theorem}
The theorem that every abelian extension of the rationals is contained in a cyclotomic field $\mathbb{Q}(\zeta_n)$; equivalently, the maximal abelian extension of $\mathbb{Q}$ is the union $\bigcup_n \mathbb{Q}(\zeta_n)$ of cyclotomic fields. By the irreducibility of the cyclotomic polynomials, the Galois group $\operatorname{Gal}(\mathbb{Q}(\zeta_n)/\mathbb{Q}) \cong (\mathbb{Z}/n\mathbb{Z})^\times$ is abelian, so the cyclotomic tower realises the whole abelian part of the absolute Galois group. The theorem is the prototype of explicit class field theory: it gives an explicit description of abelian extensions of $\mathbb{Q}$, generalised by Kronecker's Jugendtraum and the Langlands program to other number fields.

\term{Krylov Subspace}
The subspace $\mathrm{span}\{b, Ab, A^2b, \ldots\}$ used in iterative methods for linear systems such as GMRES and the conjugate gradient method.

\term{K-matrix}
In scattering theory, the K-matrix is a representation of the S-matrix that ensures unitarity. It is related to the phase shifts of the scattering process.

\term{K-theoretic Euler Characteristic}
The alternating sum of the ranks of the K-theory groups of a space, providing a generalized notion of the classical Euler characteristic.

\term{Kummer's Theorem}
A theorem giving the exponent of a prime $p$ in a binomial coefficient in terms of carries in base-$p$ addition.

\term{Kurtosis}
The fourth standardized moment of a distribution, measuring the heaviness of its tails.
\term{Opposite Category}
For a category $\mathcal{C}$, the opposite category $\mathcal{C}^{\operatorname{op}}$ has the same objects but all morphisms reversed: $\operatorname{Hom}_{\mathcal{C}^{\operatorname{op}}}(A, B) = \operatorname{Hom}_{\mathcal{C}}(B, A)$. Duality is the principle of obtaining a new theorem by reversing all arrows in a proof. The opposite of an abelian category is abelian, and $(\mathcal{C}^{\operatorname{op}})^{\operatorname{op}} \cong \mathcal{C}$.

